\subsection{Prosjektplanlegging}
\label{sec:prosjektplan}
\subsubsection{Arbeidsform og dokumentasjon}
Arbeidet gjennomføres i korte utviklingsrunder med sporbar sammenheng mellom krav, designvalg og tester. CAD-versjoner, bilder, produksjonsinnstillinger, testdata og beslutninger lagres fortløpende. Rapporten oppdateres etter hver hovedfase.

\subsubsection{Fremdriftsplan}
Prosjektet gjennomføres fra uke 39 til innlevering mandag 16. november 2026 (uke 47). Fremdriftsplanen er vist som Gantt-diagram i figur~\ref{fig:gantt}, og aktivitetene er:
\begin{enumerate}
    \item \textbf{Utrede og forstå teorien bak KERS (uke 39).} Sette seg inn i hvordan kinetisk energi kan lagres og gjenvinnes, og etablere det teoretiske grunnlaget for konseptarbeidet.
    \item \textbf{Prosjektstyring og risikovurdering (uke 40).} Planlegge gjennomføringen av prosjektet, fordele arbeid og vurdere risiko med tilhørende tiltak.
    \item \textbf{Konseptutvikling og innledende design (uke 41--43).} Utarbeide produktkrav, utvikle og vurdere konsepter, velge konsept og utforme et første design som kan testes.
    \item \textbf{Testing (uke 43--44).} Fastsette testprosedyre og teste det valgte konseptet. Testingen starter i siste uke av designfasen, slik at funn kan tas med i det videre designarbeidet.
    \item \textbf{Iterering på konsept (uke 44).} Gjennomføre avgrensede forbedringer basert på testresultatene.
    \item \textbf{Løpende rapportskriving (uke 39 til innlevering).} Rapporten skrives parallelt med de øvrige aktivitetene gjennom hele prosjektet.
    \item \textbf{Sluttføring av rapport (uke 45 til innlevering).} Ferdigstille resultater, diskusjon og konklusjoner fra testingen frem mot innleveringen 16. november.
\end{enumerate}
Aktivitetene bygger på hverandre: testing forutsetter et testbart design, og iterasjonen bygger på resultatene fra testingen.

\begin{figure}[htbp]
    \centering
    \begin{ganttchart}[
        x unit = 1.0cm,
        y unit chart = 0.75cm,
        vgrid,
        hgrid,
        title height = 1,
        title label font = \footnotesize,
        bar label font = \footnotesize,
        bar label node/.append style = {align = right, text width = 4.6cm},
        bar/.append style = {fill = NavyBlue!70, draw = NavyBlue},
        bar height = 0.5,
        milestone label font = \footnotesize\bfseries,
        milestone/.append style = {fill = BrickRed, draw = BrickRed},
        milestone label node/.append style = {align = right, text width = 4.6cm},
    ]{1}{9}
        \gantttitle{Uke 39}{1} \gantttitle{Uke 40}{1} \gantttitle{Uke 41}{1}
        \gantttitle{Uke 42}{1} \gantttitle{Uke 43}{1} \gantttitle{Uke 44}{1}
        \gantttitle{Uke 45}{1} \gantttitle{Uke 46}{1} \gantttitle{Uke 47}{1} \\
        \ganttbar{Utrede og forstå teori bak KERS}{1}{1} \\
        \ganttbar{Prosjektstyring og\\risikovurdering}{2}{2} \\
        \ganttbar{Konseptutvikling og\\innledende design}{3}{5} \\
        \ganttbar{Testing}{5}{6} \\
        \ganttbar{Iterering på konsept}{6}{6} \\
        \ganttbar[bar/.append style = {fill = NavyBlue!20}]{Løpende rapportskriving}{1}{8} \\
        \ganttbar{Sluttføring av rapport med\\konklusjoner fra testing}{7}{8} \\
        \ganttmilestone{Innlevering 16. november}{8}
    \end{ganttchart}
    \caption{Gantt-diagram for fremdriftsplanen. Løpende rapportskriving pågår gjennom hele prosjektet og er vist med lysere farge.}
    \label{fig:gantt}
\end{figure}

\subsubsection{Risikovurdering og tiltak}
Risiko vurderes som produktet av sannsynlighet og påvirkning, begge på en skala fra 1 til 5, og klassifiseres etter risikomatrisen i tabell~\ref{tab:risikomatrise}. Risikoen er delt i personskade ved det praktiske arbeidet og forhold som kan hindre gjennomføringen av prosjektet.

% Celler i risikomatrisen: \riskcell{farge}{nivå}{poeng}
\newcommand{\riskcell}[3]{\cellcolor{#1}\textbf{#2}\newline\textbf{#3}}
% Risikopoeng farget etter nivå i risikomatrisen: \riskscore{poeng}
\newcommand{\riskscore}[1]{%
    \ifnum#1>14 \cellcolor{riskextreme}\else
    \ifnum#1>7 \cellcolor{riskhigh}\else
    \ifnum#1>3 \cellcolor{riskmoderate}\else
    \cellcolor{risklow}\fi\fi\fi #1}
% Kolonner for risikotabellene; hendelsen brytes over flere linjer ved behov
\newcolumntype{E}{>{\centering\arraybackslash}m{6.4cm}}
\begin{table}[H]
    \centering
    \caption{Risikomatrise}
    \label{tab:risikomatrise}
    \small
    \renewcommand{\arraystretch}{1.5}
    \begin{tabular}{|c|>{\centering\arraybackslash}m{3.0cm}|*{5}{>{\centering\arraybackslash}m{1.6cm}|}}
        \hline
        \multicolumn{2}{|c|}{} & \multicolumn{5}{c|}{\textbf{Påvirkning}} \\
        \cline{3-7}
        \multicolumn{2}{|c|}{} & \textbf{Lav}\newline\textbf{1} & \textbf{Liten}\newline\textbf{2} & \textbf{Moderat}\newline\textbf{3} & \textbf{Høy}\newline\textbf{4} & \textbf{Alvorlig}\newline\textbf{5} \\
        \hline
        & \textbf{Svært sannsynlig}\newline\textbf{5} & \riskcell{riskmoderate}{Moderat}{5} & \riskcell{riskhigh}{Høy}{10} & \riskcell{riskextreme}{Ekstrem}{15} & \riskcell{riskextreme}{Ekstrem}{20} & \riskcell{riskextreme}{Ekstrem}{25} \\
        \cline{2-7}
        & \textbf{Sannsynlig}\newline\textbf{4} & \riskcell{riskmoderate}{Moderat}{4} & \riskcell{riskhigh}{Høy}{8} & \riskcell{riskhigh}{Høy}{12} & \riskcell{riskextreme}{Ekstrem}{16} & \riskcell{riskextreme}{Ekstrem}{20} \\
        \cline{2-7}
        & \textbf{Nøytral}\newline\textbf{3} & \riskcell{risklow}{Lav}{3} & \riskcell{riskmoderate}{Moderat}{6} & \riskcell{riskhigh}{Høy}{9} & \riskcell{riskhigh}{Høy}{12} & \riskcell{riskextreme}{Ekstrem}{15} \\
        \cline{2-7}
        & \textbf{Lite sannsynlig}\newline\textbf{2} & \riskcell{risklow}{Lav}{2} & \riskcell{riskmoderate}{Moderat}{4} & \riskcell{riskmoderate}{Moderat}{6} & \riskcell{riskhigh}{Høy}{8} & \riskcell{riskhigh}{Høy}{10} \\
        \cline{2-7}
        \multirow{-10}{*}{\rotatebox[origin=c]{90}{\textbf{Sannsynlighet}}} & \textbf{Sjelden}\newline\textbf{1} & \riskcell{risklow}{Lav}{1} & \riskcell{risklow}{Lav}{2} & \riskcell{risklow}{Lav}{3} & \riskcell{riskmoderate}{Moderat}{4} & \riskcell{riskmoderate}{Moderat}{5} \\
        \hline
    \end{tabular}
\end{table}

\paragraph{Personskade}
Tabell~\ref{tab:risiko-arbeid} viser HMS-relatert risiko knyttet til det praktiske arbeidet.

\begin{table}[H]
    \centering
    \caption{Vurdering av potensielle skadefaktorer}
    \label{tab:risiko-arbeid}
    \small
    \begin{tabular}{|c|E|c|c|c|}
        \hline
        nr & \textbf{Potensielle hendelser} & \textbf{Sannsynlighet} & \textbf{Påvirkning} & \textbf{Risikopoeng} \\
        \hline
        \textbf{1.} & Skader fra bruk av dreiebenk & 3 & 4 & \riskscore{12} \\
        \hline
        \textbf{2.} & Brannskade ved sveising & 2 & 5 & \riskscore{10} \\
        \hline
        \textbf{3.} & Brannskade av 3D-printer & 1 & 2 & \riskscore{2} \\
        \hline
        \textbf{4.} & Klemmer finger i bilmekanisme & 1 & 3 & \riskscore{3} \\
        \hline
        \textbf{5.} & Skade av verktøy (Ubåten) & 2 & 3 & \riskscore{6} \\
        \hline
    \end{tabular}
\end{table}

Tiltak for å redusere risiko for personskade, nummerert etter hendelsene i tabell~\ref{tab:risiko-arbeid}:
\begin{enumerate}
    \item \textbf{Dreiebenk:} Få opplæring før bruk og følg verkstedets instrukser. Bruk vernebriller, unngå løse klær, smykker og løst hår, og forlat aldri maskinen mens den går.
    \item \textbf{Sveising:} Følg alle instrukser knyttet til bruk av sveiseapparat. Blant annet: ikke bruk plagg med lange ermer, bruk påbudt sikkerhetsutstyr til enhver tid og ha brannslukningsapparat tilgjengelig.
    \item \textbf{3D-printer:} Ikke ta på dysen eller printplaten før printen er ferdig og har kjølt seg ned.
    \item \textbf{Bilmekanismen:} Kontroller innfesting før opptrekk, og hold hender og løse gjenstander unna roterende deler.
    \item \textbf{Verktøy i Ubåten:} Utvis sunn varsomhet, sett av god tid til oppgavene som skal utføres, bruk nødvendig verneutstyr og følg verkstedets sikkerhetsrutiner.
\end{enumerate}

\paragraph{Gjennomføring av prosjektet}
Tabell~\ref{tab:risiko-gjennomforing} viser hendelser som kan forsinke eller hindre prosjektet.

\begin{table}[H]
    \centering
    \caption{Potensielle hendelser under gjennomføring av prosjektet}
    \label{tab:risiko-gjennomforing}
    \small
    \begin{tabular}{|c|E|c|c|c|}
        \hline
        nr & \textbf{Potensielle hendelser} & \textbf{Sannsynlighet} & \textbf{Påvirkning} & \textbf{Risikopoeng} \\
        \hline
        \textbf{1.} & 3D-print feiler & 4 & 2 & \riskscore{8} \\
        \hline
        \textbf{2.} & Gruppemedlem blir syk & 1 & 3 & \riskscore{3} \\
        \hline
        \textbf{3.} & Konsept fungerer ikke & 3 & 3 & \riskscore{9} \\
        \hline
        \textbf{4.} & Deler blir skadet & 4 & 3 & \riskscore{12} \\
        \hline
        \textbf{5.} & Utdelte komponenter har avvikende dimensjoner eller skjevheter som må rettes & 2 & 3 & \riskscore{6} \\
        \hline
    \end{tabular}
\end{table}
Påvirkningen for hendelse 4 i tabell~\ref{tab:risiko-gjennomforing} avhenger av hvilken del som blir skadet. 3D-printede deler er lette å bytte ut, mens for eksempel svinghjulet er vanskeligere å erstatte.

Tiltak for å redusere påvirkning på prosjektets fremgang, nummerert etter hendelsene i tabell~\ref{tab:risiko-gjennomforing}:
\begin{enumerate}
    \item \textbf{Feilet 3D-print:} Rengjør printplaten med sprit før utskrift og følg med på minst de to første lagene. Start utskrift av kritiske deler tidlig, slik at det er tid til å printe på nytt.
    \item \textbf{Sykdom:} Gruppen består bare av ett medlem. Ved langvarig sykdom må evt. veileder kontaktes. 
    \item \textbf{Konseptet fungerer ikke:} Test delsystemer tidlig, allerede mens designet utvikles, og unngå store, drastiske designendringer sent i prosjektet.
    \item \textbf{Skadde deler:} Ha reservedeler av 3D-printede komponenter og ferdige printfiler klare, slik at skadde deler kan erstattes raskt. Deler som er vanskelige å erstatte, som svinghjulet, håndteres og lagres forsiktig.
    \item \textbf{Avvikende utdelte komponenter:} Mål opp utdelte komponenter ved mottak, kontroller kritiske mål og test delsystemer før full montering, slik at skjevheter og dårlig passform oppdages tidlig.
\end{enumerate}
Generelle tiltak som ikke knytter seg til én enkelt hendelse:
\begin{itemize}
    \item Leveringstid og verkstedtilgang: avklar kritiske deler og bestill eller reserver tid tidlig.
    \item Varierende opptrekk og underlag: standardiser testprosedyren og registrer avvik.
\end{itemize}
